% 使用 XeLaTeX 编译。
\documentclass[UTF8,a4paper,11pt,fontset=fandol]{ctexart}
\usepackage[margin=2.5cm]{geometry}
\usepackage{amsmath,amssymb}
\usepackage{booktabs,tabularx,array}
\usepackage[hidelinks]{hyperref}
\setlength{\parskip}{0.4em}
\title{无人机 MAPPO 决策替换方案概要}
\author{Teamed-Planning 项目}
\date{2026年10月12日}

\begin{document}
\maketitle

\section{方案目标与范围}
采用多智能体近端策略优化（MAPPO）设计无人机航点选择算法，实现多架无人机同时决策、同时移动、同时观测。优化整体地图置信度，并使无人车附近、任务目标周围以及无人车与目标之间走廊的有效置信度增益占比较高，同时保留其他区域的探索能力。

强化学习的决策主体是无人机，无人车导航算法暂不纳入优化。\textbf{目标点不进入部署时的无人机策略输入}；目标相关区域仅由训练端定义并评价。没有允许接收的目标相关提示时，策略不能保证在任意隐藏目标下定向提高对应走廊置信度，相关贡献只能作为学习目标与独立测试指标。

本文件仅修改方案，不代表算法已经实现或模型已经训练。按用户明确的信息边界，每架无人机仅维护由自身观测更新的置信图，地图上传到无人车端进行融合；无人机之间不交换或融合地图。地图中继若存在，只转发数据，不将载荷作为中继无人机的地图或策略输入。

\subsection{原设计与当前代码的差异}
当前仓库的 \texttt{drone\_main.py} 含有读取邻机图并融合的分支，任务适配器还读取目标和规划路径；\texttt{ros\_drone\_node.py} 也接收这些消息。因此当前代码与用户描述的原始信息边界不完全一致。方案沿用原接口形式及可用的局部字段，不新增全局视角，但明确排除目标、完整规划路径和邻机地图。后续实现必须先统一该边界，不能宣称当前代码已经满足要求。原 C-ADMM 本身依赖目标；保留目标输入的原基线与本方案存在信息差异，应单独报告，不能作为同信息条件下的公平比较。

\section{训练架构与信息边界}
采用集中训练、分散执行。所有无人机共享 Actor 参数，部署时各机根据自己的本地信息独立决策；集中式 Critic 只在训练时使用。

\begin{table}[htbp]
\centering
\begin{tabularx}{\textwidth}{@{}p{2.4cm}X@{}}
\toprule
模块 & 输入与限制 \\
\midrule
Actor & 本机三维位置、本机感知构建的独立置信图及历史、允许接收的邻机位置与拟定航点、已收到的无人车位置、时间及现有 ID、通信距离等配置。\textbf{不输入目标点、完整规划路径、邻机地图或无人车融合地图}。\\
Critic & 汇总同一训练时间步各无人机合法的本地观测，使用有效性掩码处理缺失信息。不读取全局真实地图、未观测障碍物或其他特权真值。\\
奖励模块 & 在训练端使用无人车实际收到的有效证据与融合结果，按来源去重；目标坐标仅用于定义任务区域。融合地图、观测账本及目标信息均不回传为 Actor 输入。\\
\bottomrule
\end{tabularx}
\end{table}

集中式 Critic 扩大的是训练端对合法本地观测的汇总范围。执行端不增加必需消息字段，邻机状态与航点只在原接口允许接收时使用。速度、剩余能量、链路质量及上传确认均不能被默认视为已有输入；运动趋势可由收到的位置历史估计，并应考虑误差。

仿真中的全部机器人真实位置和速度、完整障碍物列表、未来轨迹等字段不得进入 Actor、Critic 或奖励计算。环境真值仅用于物理演化和独立评估。训练观测中的图属性应与部署接口一致，避免 Python 仿真保留了 ROS 消息未传输的属性而造成信息差异。

\section{决策接口与动作空间}
保持现有接口：
\begin{quote}
\texttt{compute\_next\_target(current\_position, observation\_data)}\\
返回目标航点 $(x,y,z)$ 或无目标。
\end{quote}

新增 RL 航点选择器，以配置切换当前 C-ADMM 和 MAPPO。初始沿用默认每秒一次的调度周期，底层控制继续按现有频率运行。

动作是在可变候选航点集合中选择一个，另提供保持当前位置的选项：
\[
 a_i^t\in\{1,\ldots,K_i^t,\mathrm{hold}\},\qquad
 w_i^t=(x_{a_i^t},y_{a_i^t},z_i^t).
\]
候选航点保持本机当前飞行高度。复用仅依赖本机图的低观测次数节点、断连区域中点及局部前沿逻辑，并加入其他方向的探索候选。初始约预留 25\% 的候选名额用于自由探索；该比例不等于强制执行时间比例。

网络对每个候选点计算分数，使用动作掩码屏蔽空位和不合法动作，不把固定动作编号绑定到固定空间位置。候选点及其特征仅由已有信息计算，不预先读取未知区域的真实收益。

由于目标和完整规划路径不可见，部署候选生成不能沿用依赖这些字段的路径中点、指向目标的前沿或直线回退。改为依据本机低置信度节点、局部前沿、其他方向的探索点以及收到的无人车位置生成候选。策略可使用历史记忆；历史记忆不能恢复从未提供的隐藏目标信息。

\subsection{同时执行与局部协同}
每个决策周期各机从同一轮可用局部信息独立提出动作，通过已有邻机位置与航点进行局部可行性检查。仿真先收集整队动作，再统一推进同一个时间步；部署时各机持续跟踪自己的航点，不等待其他无人机到达。执行期间可继续使用上一段有效动作，避免通信等待演变为轮流飞行。这里的并行指实际多机同时移动，不是多个训练环境并行。

当前二维仿真已经收集各机动作后按同一时间步推进；控制器在程序中依次计算不能证明无人机轮流移动。新方案应减少动作冲突和停滞，并通过同时有效移动比例评价，而不预设原实现为串行飞行。

\section{探索区域与最小奖励}
\subsection{空间优先级}
\begin{enumerate}
  \item 无人车附近：以收到的无人车当前位置为中心，半径为 $R_c$。
  \item 目标点周围：训练端以任务目标为中心，半径为 $R_g$。
  \item 无人车到目标之间：训练端以无人车当前位置与目标连线段为中心，半宽为 $B$ 的带状区域，允许向两侧探索绕行空间。
  \item 自由探索区域：上述重点区域之外的可探索空间。
\end{enumerate}
重点区域取并集，仅用于训练奖励与独立评估，不作为无人机输入。半径、带宽和权重可配置。执行端只能根据自身置信图和收到的无人车位置形成可观测偏好；若无人车运动趋势可从历史估计，可优先探索其前方，但这不等价于知道目标走廊。

\subsection{统一团队奖励}
用不重叠的重点区域 $\Omega_T$ 和其他区域 $\Omega_E$ 分别统计有效置信度增益。设 $C_t(q)$ 为无人车实际收到证据后对位置 $q$ 的置信度估计。每周期采用动作执行前确定的同一组区域，定义
\[
 G_X^t=\sum_{q\in\Omega_X^t} A(q)\bigl[C_{t+1}(q)-C_t(q)\bigr],\qquad X\in\{T,E\}.
\]
整体收益为 $G_T^t+G_E^t$，重点收益另有加权：
\begin{equation}
 r_t=\alpha\frac{G_T^t+G_E^t}{Z_C}
 +\beta_t\frac{G_T^t}{Z_C}
 +\eta_t\frac{A_{E,\mathrm{new}}^t}{Z_A}
 -\lambda\frac{D_t}{Z_D}-c\Delta t,
 \label{eq:reward}
\end{equation}
其中 $A_{E,\mathrm{new}}^t$ 是实际送达且有效的外部首次观测面积，$D_t$ 是团队飞行距离，作为未引入电量传感器时的运动代价代理。$Z_C,Z_A,Z_D$ 为固定归一化尺度；外部新增探索项设上限，避免压过置信度提升。重点与其他区域互不重叠，整体项和重点项的双重奖励明确表示任务偏好。

障碍物与可通行区域的有效确认均有价值。重复观测只要降低不确定性或消除冲突，仍可产生收益；重复上传或循环融合同一证据不能加分。置信度不能仅由观测次数无限增长，需要质量加权、来源去重和独立校准；相互矛盾的证据导致置信度下降时保留负增益，不通过截断负值鼓励过度自信。随着机器人移动改变区域定义，不会重置首次观测账本。

\subsection{任务贡献占比与探索平衡}
在固定时长滑动窗口内分别累计重点区域与其他区域的正向有效增益 $G_T^+,G_E^+$，定义
\[
 \rho_T=\frac{G_T^+}{G_T^++G_E^++\varepsilon}.
\]
初始可将 $0.60\leq\rho_T\leq0.75$ 作为待验证的软目标。占比过低时缓慢增加 $\beta_t$，过高时增加 $\eta_t$，并对权重设置上下界；总增益接近零时不更新。置信度下降与冲突另行评价，不能被该正向占比掩盖。重点区域达到质量阈值后逐渐降低额外奖励，恢复外部探索。通信可达性与隐藏目标使该占比不能成为逐步执行保证，实验应报告可达区域和未达标原因。

观测账本由实际到达的有效观测构建，不能把未观测但被当前程序当作空闲的位置视为已观测，也不能仅凭相机预计覆盖范围给分。模型缺失时生成的全零回退占用图不应作为有效探索证据。

暂不加入无人车路径改善、接近目标或导航进度奖励，无人车导航不参与此次优化。时间代价在固定长度回合中不能单独区分不同轨迹，因此以有效地图增益、探索收益及运动代价共同评价。回合长度及终止规则固定，避免通过异常提前终止逃避代价。

\section{约束与回退}
通信距离、无人机间距和高度限制通过动作筛选及控制层处理，第一版不分别加入奖励。筛选只使用已接收的邻机位置与航点，预测邻机也在移动时的短时轨迹，并检查全过程三维距离。初始采用预配置中继链角色，缺失邻居不能自动从链中删除并宣称仍然连通。直接上传模式与允许中继转发模式分别配置。

没有可行动作或推理异常时，使用明确的悬停或恢复通信规则，并记录回退次数。训练和执行采用一致的动作筛选方式。局部信息可能延迟，动作筛选不能保证飞行全过程绝对连通或无碰撞，应通过独立指标验证。

当前占用图主要描述二维地面通行性，不能直接视为无人机飞行高度上的三维障碍地图。本方案不新增三维感知能力。

\section{训练、验证与验收}
\begin{enumerate}
  \item 建立输入白名单和训练环境适配层，隔离仿真真值，统一训练与部署的数据字段。
  \item 在现有 Python 仿真中训练，从单机简单场景扩展到多机、复杂地图和通信延迟；随机化地图与任务，分离训练、验证和测试场景。
  \item 按联合仿真时间步收集样本，训练共享 Actor 与集中式 Critic；每个决策步包含原有控制和地图更新过程。
  \item 在 XTDrone 中验证图像分割误差、通信延迟和飞行动力学对策略的影响。
\end{enumerate}

主要评估整体地图置信度增益与校准误差、重点区域增益占比、外部有效探索面积、无效重复观测、同时有效移动比例、飞行距离、通信中断、间距违约、碰撞、回退频率及决策耗时。无人车导航指标不作为本阶段验收要求。在相同局部信息、候选集合和控制条件下比较贪心探索基线；原始依赖目标及邻机融合的 C-ADMM 单独列为具有不同信息条件的参考，不能把信息差异归因于强化学习效果。使用多个随机种子和独立测试地图、目标验证。

采用固定决策预算或明确的任务终止条件，不按策略自行申报的探索完成状态结束回合。通过消融实验分别验证整体增益、任务偏好和外部探索奖励的作用，并评价隐藏目标下可辨识性的限制。MAPPO 与奖励设计均不构成收敛或全局最优保证。

\section{拟新增模块与当前状态}
后续实现拟新增观测编码器、MAPPO 航点选择器、训练环境适配器、实际观测去重账本和奖励模块，并提供算法切换配置。主要接入位置为 \texttt{drone\_navigator.py} 的决策接口和 \texttt{drone\_main.py} 的调度调用。

当前阶段仅形成方案文档，尚未修改无人机运行算法。重点区域尺寸、时间代价与归一化尺度在实现阶段依据当前地图分辨率、相机覆盖范围和回合预算确定。

\end{document}
